\chapter{Strain-Rate Dependent Mechanical Behavior of HANI Constructs}\label{chap:6}


\section{Dynamic Mechanical Loading and Strain Rate Sensitivity of the Meniscus}

During activities such as walking, running, and squatting, the knee joint experiences complex dynamic loading characterized by rapid fluctuations in compressive force. Within this environment, the meniscus, a fibrocartilaginous tissue interposed between the femoral condyles and tibial plateau, plays a vital role in distributing loads, enhancing joint congruity, and protecting articular cartilage surfaces \cite{thambyah2006mechanicalproperties}. Critical to the meniscus's functionality is its ability to accommodate a wide range of compressive strain rates, with localized axial strains reaching up to 15\% under normal physiological loading conditions \cite{morejon2020compressiveproperties}.

The mechanical response of the meniscus varies significantly with loading rate, demonstrating a sophisticated strain-rate dependent behavior that optimizes load distribution across different activities \cite{danso2014comparisonnonlinear}. During slow loading (e.g., standing), the tissue exhibits fluid-dominated viscoelastic behavior, allowing for gradual stress redistribution. In contrast, rapid loading scenarios (e.g., jumping) trigger a transition to fiber-dominated mechanics, where aligned collagen networks provide immediate structural resistance critical for joint protection \cite{morejon2023tensileenergy, norberg2021viscoelasticequilibrium}.

The mechanical evaluation of tissue-engineered meniscal replacements must therefore account not only for peak mechanical properties but also for their behavior under dynamically applied loads. Previous investigations have largely focused on quasi-static or discrete stress-relaxation protocols to characterize meniscal constructs. While informative, these approaches fail to replicate the continuous and variable nature of in vivo strain, limiting their predictive value for dynamic joint performance.

In this study, the mechanical behavior of hydrogel-PCL composites, formulated using the Hybrid Hydrogels Augmented via Additive Network Integration (HANI) method, is evaluated under continuous compressive loading at physiologically relevant strain rates. Through this approach, we aim to capture the influence of loading rate on construct stiffness and to assess the ability of different PCL network designs to withstand dynamic compressive forces while maintaining mechanical integrity.

\section{Hypothesized Mechanical Response}
Fiber orientation and network connectivity are hypothesized to significantly influence scaffold strain-stiffening behavior under dynamic compressive loading. The 3D Aligned network, with its helically oriented fibers, is expected to exhibit the greatest strain-rate sensitivity due to progressive fiber recruitment during compression. In contrast, the Gyroid network's isotropic architecture should result in minimal rate-dependent behavior. The X Pattern network, featuring intermediate structural anisotropy, is anticipated to display moderate rate sensitivity, bridging the mechanical response gap between the other two designs.

\section{Design Strategy for HANI Constructs Under Dynamic Compression}

The mechanical behavior of native meniscal tissue is governed by its hierarchical poro-viscoelastic structure, wherein the interplay between solid matrix elasticity, fiber orientation, and interstitial fluid flow determines the tissue's load-bearing and energy-dissipating capabilities \cite{kleinhans2018hydraulicpermeability, morejon2020compressiveproperties}. At low strain rates, the meniscus exhibits fluid-mediated viscoelastic relaxation, while at higher rates, the solid fiber network increasingly dominates the mechanical response \cite{abdelgaied2015comparisonbiomechanical}.

Three distinct PCL network designs and one thermomolded PCL structure were fabricated to systematically explore the relationship between architectural anisotropy and strain-rate dependent mechanical behavior. By integrating these designs within a common hydrogel base and subjecting the resulting constructs to confined compression at multiple strain rates, this study aims to identify key structure-function relationships governing dynamic load transmission and mechanical resilience in engineered meniscal replacements.

The selection of distinct PCL network designs was guided by biomimetic principles established in recent tissue engineering approaches \cite{szojka2017biomimetic3d}. The 3D Aligned network, featuring helically oriented fibers, mimics circumferential collagen architecture. The Gyroid network offers isotropic stress distribution. The X Pattern network, with intersecting fibers at ±30°, provides intermediate anisotropy. This strategic selection allows systematic exploration of how directional reinforcement and architectural complexity modulate strain-rate dependent mechanical behavior.

\section{Experimental Design Flow}

To rigorously evaluate the strain-rate dependent mechanical behavior of HANI constructs, a structured experimental workflow was developed. This approach was designed to systematically isolate the effects of PCL network architecture while maintaining consistent hydrogel composition and scaffold geometry across all test groups. The experimental protocol was divided into five key stages: hydrogel fabrication, PCL network integration, confined compression testing, mechanical analysis, and statistical evaluation.

\subsection{Hydrogel Fabrication}

The hydrogel matrix was fabricated using Type A gelatin crosslinked with glutaraldehyde (GTA) at a standardized 200:1 mass ratio (GelA:GTA). This formulation was selected based on prior studies demonstrating its ability to replicate physiologically relevant water content and viscoelastic behavior characteristic of native meniscal tissue. All hydrogels were molded using a custom-designed 3D thermomold, featuring a geometry with a 6 mm diameter top surface elevated 4 mm above a 12 mm diameter circular base. This shape was specifically chosen to replicate the dome-like morphology of the meniscus and to impose geometric boundary conditions that simulate in vivo mechanical constraints. Fabrication processes were optimized to ensure homogeneous crosslinking and reproducibility across batches.

\subsection{PCL Network Integration}

Three distinct PCL network architectures were embedded within the gelatin hydrogel during the molding process. The network designs were selected to explore a spectrum of architectural anisotropies:

\begin{itemize}
    \item \textbf{3D Aligned Network:} Comprising helically oriented fibers with 1 mm spacing and 5 mm radial tie fibers, mimicking the circumferential collagen fiber arrangement of the native meniscus.
    \item \textbf{Gyroid Network:} A triply periodic minimal surface (TPMS) structure with 1.8 mm feature spacing, designed to provide isotropic mechanical support.
    \item \textbf{X Pattern Network:} A crisscross lattice with fibers rotated at +30$^\circ$ and -30$^\circ$ relative to the vertical axis, spaced 2.1 mm apart, offering intermediate anisotropy.
\end{itemize}

Integration of the PCL networks was performed within the mold during hydrogel casting to ensure complete encapsulation and uniform embedding of the reinforcement structures.

\subsection{Confined Compression Testing}

Mechanical testing was conducted using a custom-built confined compression chamber designed to apply uniaxial compressive loads under controlled conditions. Each sample was subjected to continuous compressive loading from 5\% to 25\% strain, with no intermediate hold periods, to mimic dynamic physiological loading scenarios. Three strain rates were selected for testing:

\begin{itemize}
    \item 0.1\%/s (slow, simulating quasi-static loading such as standing),
    \item 1\%/s (moderate, simulating typical dynamic activities such as walking),
    \item 5\%/s (fast, simulating high-impact loading such as jumping).
\end{itemize}

This range of strain rates was chosen to encompass the full spectrum of physiological loading conditions experienced by the meniscus in vivo, ensuring the relevance of the findings.

\subsection{Mechanical Analysis}

Force and displacement data were recorded continuously throughout the testing process. Engineering stress and strain were calculated for each sample, and the confined Young's modulus ($E_{\text{confined}}$) was determined from the linear portion of each stress-strain curve, typically between 10\% and 20\% strain. This modulus served as the primary metric for evaluating the stiffness and load-bearing capacity of the constructs across different strain rates and network designs.

\section{Objective}

The primary objective of this study is to investigate the confined compressive mechanical behavior of HANI constructs under continuous dynamic loading conditions, with a specific focus on the influence of strain rate and internal PCL network architecture. Unlike conventional mechanical testing protocols that rely on static or discrete hold periods to characterize scaffold behavior, this study employs a continuous compression protocol designed to more closely replicate the dynamic, variable strain conditions experienced by the meniscus in vivo.

All constructs were fabricated using a standardized gelatin hydrogel formulation (GelA:GTA mass ratio of 200:1) and molded into a consistent three-dimensional geometry featuring a 6 mm diameter top surface elevated 4 mm above a 12 mm diameter circular base. This controlled geometry ensures that any observed differences in mechanical response are attributable solely to variations in PCL network architecture, thereby isolating internal structure as the primary independent variable.

Specifically, the study aims to:

\begin{itemize}
    \item Quantify the confined compressive stiffness ($E_{\text{confined}}$) of HANI constructs across multiple physiologically relevant strain rates (0.1\%/s, 1\%/s, and 5\%/s).
    \item Assess the role of PCL network patterning in modulating strain-rate sensitivity and mechanical reinforcement behavior.
    \item Identify scaffold architectures that best replicate the dynamic load-bearing performance of the native meniscus, with a focus on optimizing mechanical robustness and compliance.
\end{itemize}

Dynamic performance is evaluated according to several key criteria:

\begin{itemize}
    \item Progressive increases in confined modulus with increasing strain rate, indicating effective strain-rate sensitivity.
    \item Maintenance of mechanical compliance and avoidance of brittle failure across all testing conditions.
    \item Statistical significance of strain-rate and network architecture effects on scaffold stiffness and load transfer behavior.
\end{itemize}

Through this systematic investigation, the study seeks to define critical structure-function relationships that inform the rational design of meniscal replacements capable of withstanding the complex biomechanical demands of the knee joint.



\section{Objective}

The objective of this study is to investigate the strain-rate dependent confined compressive mechanical behavior of composite HANI constructs featuring distinct PCL network designs. Unlike traditional stress-relaxation experiments that focus on discrete hold periods at fixed strains, this work employs continuous compressive loading from 5\% to 25\% strain. This protocol enables evaluation of mechanical response under dynamic strain conditions that more accurately mimic the continuous, variable loading encountered by the meniscus in vivo.

A critical feature of the experimental design is the isolation of PCL reinforcement pattern as the primary independent variable. All constructs were fabricated using a standardized GelA:GTA hydrogel formulation (200:1 mass ratio) and molded into an identical 3D geometry featuring a 6 mm diameter top surface elevated 4 mm above a 12 mm diameter base. This configuration creates a domed structure that imposes geometric boundary conditions relevant to meniscal morphology while ensuring that any differences in mechanical behavior arise from network architecture rather than from macroscopic shape or size differences.

Specifically, this study aims to:
\begin{itemize}
    \item Quantify the confined compressive stiffness ($E_{\text{confined}}$) of HANI constructs across multiple strain rates representative of physiologic activities.
    \item Determine the influence of PCL network pattern on strain-rate sensitivity and mechanical load transfer mechanisms.
    \item Identify structural reinforcement strategies that best replicate the dynamic mechanical behavior required for successful meniscal tissue engineering applications.
\end{itemize}

By systematically varying PCL network design while maintaining consistent hydrogel properties and construct geometry, this work seeks to define structure-function relationships critical to optimizing dynamic performance of engineered meniscal replacements.

Dynamic performance will be evaluated according to the following criteria:
\begin{itemize}
    \item Progressive increases in confined modulus with increasing strain rate, indicating effective strain-rate sensitivity
    \item Maintenance of construct compliance across all testing conditions
    \item Absence of brittle failure or permanent deformation during high-rate loading
    \item Statistical significance of strain-rate effects on mechanical properties
\end{itemize}

\section{PCL Network Designs Evaluated}

Three distinct PCL reinforcement designs were fabricated, integrated within hydrogel matrices, and evaluated to systematically assess the role of internal network architecture on mechanical performance under varying compressive strain rates. All constructs were molded onto an identical three-dimensional scaffold geometry, featuring a 6 mm diameter top surface elevated 4 mm above a 12 mm diameter circular base. This thermomolded geometry imposes a biomimetic dome-like curvature relevant to meniscal anatomy and ensures that macroscopic shape and size are consistent across all test groups.

The PCL network designs evaluated were:

\begin{itemize}
    \item \textbf{3D Aligned Network:}  
    A spring-like, helically oriented arrangement of PCL struts fabricated with 1 mm aligned fiber spacing. This anisotropic design mimics the circumferential collagen fiber organization of native meniscal tissue, promoting fiber recruitment and strain-stiffening behavior under dynamic compressive loading.

    \item \textbf{Gyroid Network:}  
    A triply periodic minimal surface (TPMS) architecture, offering a continuous and isotropic load distribution. Fabricated via fused deposition modeling (FDM), this structure lacks preferential fiber orientation and is expected to exhibit relatively rate-independent mechanical behavior across strain rates.

    \item \textbf{X Pattern Network:}  
    A crisscrossed lattice of orthogonally intersecting PCL struts, providing moderate anisotropy through two principal fiber axes. This architecture enables partial fiber alignment and directional load transfer, and is hypothesized to display intermediate mechanical behavior relative to fully anisotropic and isotropic networks.
\end{itemize}

Each PCL network was embedded within a standardized hydrogel matrix (GelA:GTA mass ratio 200:1), molded into the same three-dimensional configuration, and tested under confined compression at multiple strain rates. This systematic framework isolates the influence of internal network architecture on strain-rate dependent mechanical behavior, providing critical insights for the design of dynamic load-bearing meniscal replacements.

\begin{table}[H]
\centering
\renewcommand{\arraystretch}{1.3}
\setlength{\tabcolsep}{12pt}
\caption{\textbf{Experimental Design Parameters for HANI Construct Testing.} Comprehensive overview of the key experimental parameters used to evaluate strain rate-dependent mechanical properties of hydrogel-PCL composite constructs, including material specifications, geometric constraints, testing conditions, and analysis methods.}
\label{tab:test_design}
\begin{tabular}{|p{4cm}|p{10cm}|}
\hline
\rowcolor{gray!20}
\textbf{Component} & \textbf{Specification} \\ \hline
Hydrogel Base & Gelatin A crosslinked with Glutaraldehyde (GelA:GTA mass ratio 200:1) \\ \hline
Mold Geometry & 6 mm diameter top surface elevated 4 mm above 12 mm diameter circular base \\ \hline
PCL Network Patterns & Aligned (1 mm fiber spacing with 5 mm radial ties); Gyroid (1.8 mm spacing); X Pattern (+30$^\circ$ and -30$^\circ$ rotation, 2.1 mm spacing) \\ \hline
Construct Diameter & 12.7 mm (after molding) \\ \hline
Construct Height & 4 mm \\ \hline
Testing Protocol & Confined compression testing \\ \hline
Strain Range & 5\% to 25\% compressive strain (continuous) \\ \hline
Strain Rates Tested & 0.1\%/s (slow), 1\%/s (moderate), 5\%/s (fast) \\ \hline
Primary Mechanical Metric & Confined Young's Modulus ($E_{\text{confined}}$) extracted from stress-strain curves \\ \hline
Sample Size per Group & \(n = 5\) \\ \hline
Analysis Method & Two-way ANOVA (Factors: PCL Pattern, Strain Rate) \\ \hline
\end{tabular}
\end{table}

\begin{figure}[h!]
\centering
\begin{tikzpicture}[
    node distance=2cm,
    every node/.style={
        draw,
        align=center,
        rectangle,
        rounded corners,
        minimum width=4.5cm,
        minimum height=1.2cm,
        fill=gray!10,
        font=\sffamily\small
    },
    every path/.style={
        ->,
        >=Stealth,
        thick,
        draw=blue!70
    }
]

% Nodes
\node (start) {Hydrogel Fabrication\\ GelA:GTA (200:1)\\ 3D Thermomold Geometry};
\node (pcl) [below of=start] {PCL Network Integration\\ Aligned, Gyroid, X Pattern};
\node (test) [below of=pcl] {Confined Compression Testing\\ 5\%--25\% strain\\ 0.1\%/s, 1\%/s, 5\%/s};
\node (analysis) [below of=test] {Mechanical Analysis\\ Calculate $E_{\text{confined}}$};
\node (stats) [below of=analysis] {Statistical Evaluation\\ Two-way ANOVA};

% Arrows
\draw (start) -- (pcl);
\draw (pcl) -- (test);
\draw (test) -- (analysis);
\draw (analysis) -- (stats);

\end{tikzpicture}
\caption{Experimental workflow for strain rate-dependent mechanical testing of HANI constructs.}
\label{fig:experimental_flowchart}
\end{figure}
